\section{Conclusion and Recommendations}

\subsection{Conclusion}
Among the evaluated models, the Support Vector Machine (SVM) achieved the highest overall performance, demonstrating superior classification capability compared to k-Nearest Neighbors (k-NN), Random Forest, and Naive Bayes. The SVM model obtained the highest accuracy and F1-score, indicating its effectiveness in learning discriminative patterns from the extracted image features and handling variations in hand gestures. The confusion matrix analysis further confirmed the reliability of the model, as most predictions were correctly classified with minimal misclassification across the ASL alphabet classes.

The results also showed that Random Forest and k-NN produced competitive performance, suggesting that these models are also capable of recognizing ASL gestures with high accuracy. In contrast, Naive Bayes exhibited significantly lower performance, indicating limitations in modeling the complex and non-linear relationships present in hand gesture image data.

Overall, the findings of the study demonstrate that machine learning techniques, particularly SVM, can effectively support the automated recognition of American Sign Language. The study highlights the potential of machine learning-based ASL recognition systems in improving accessibility and communication for individuals who rely on sign language.

\subsection{Recommendations}

Based on the findings of the study, several recommendations are proposed for future research and improvement of American Sign Language (ASL) recognition systems. Future studies may utilize larger and more diverse ASL datasets to improve the robustness and generalization capability of machine learning models across different hand shapes, lighting conditions, and backgrounds. Additional preprocessing and feature extraction techniques may also be explored to further enhance classification accuracy and minimize misclassification between visually similar gestures.

Moreover, future researchers may investigate the use of advanced deep learning architectures such as Convolutional Neural Networks (CNN), EfficientNet, ResNet, and Vision Transformers to compare their performance with traditional machine learning models. Since this study focused primarily on static ASL alphabet recognition, future work may also consider dynamic sign recognition involving motion-based gestures and continuous sign language translation to create more comprehensive communication systems.

The development of a real-time ASL recognition application using webcams, mobile devices, or embedded systems is also recommended to evaluate the practical applicability of the proposed approach in real-world environments. Furthermore, future studies may include additional evaluation metrics such as computational efficiency, inference time, and robustness under varying environmental conditions. Integrating ASL recognition systems with speech synthesis or text generation technologies may also further improve accessibility and communication support for deaf and hard-of-hearing individuals.